\section{Related Work}
\label{sec:related}

\para{Residual stream analysis and mechanistic interpretability.}
The residual stream---the sum of all previous layer outputs---serves as the primary information conduit in transformers. The \emph{logit lens}~\citep{nostalgebraist2020logitlens} and its refinement, the \emph{tuned lens}~\citep{belrose2023tuned}, decode intermediate hidden states into vocabulary distributions, revealing how predictions evolve layer by layer. \citet{lawson2025mlsae} extend this with multi-layer sparse autoencoders (MLSAEs) that decompose the residual stream into interpretable features across all layers simultaneously.
Circuit discovery methods such as ACDC~\citep{conmy2023automated} and edge attribution patching~\citep{syed2023attribution} identify which components drive specific behaviors. \citet{zhao2024residual} show that knowledge conflicts between parametric and contextual information create detectable patterns in the residual stream. \citet{nanda2023emergent} demonstrate that transformers learn linear world-model representations in the residual stream. Our work extends this line of research by studying how the residual stream responds to systematic input perturbations rather than natural language variation.

\para{Token permutation and cipher robustness.}
\citet{sinha2021masked} show that masked language models pre-trained on randomly shuffled text retain 95--98\% of downstream performance on most tasks, demonstrating that distributional co-occurrence statistics---not word order---drive much of pre-training success. \citet{ravishankar2022word} challenge this finding by showing shuffled-text models still encode some positional information through architectural inductive biases. \citet{haviv2022positional} further demonstrate that transformers without explicit positional encodings learn position from causal attention masks alone. These results suggest that transformers have inherent structural biases that persist regardless of input ordering, which is relevant to understanding how the residual stream processes bijectively permuted tokens.

On the cipher side, \citet{yuan2023gpt4cipher} show that GPT-4 can interpret prompts encoded with Caesar ciphers, Morse code, and ASCII encoding---even without explicit demonstrations---indicating that large models develop latent cipher-processing capabilities during pre-training. \citet{li2025cipherbank} provide a systematic benchmark finding that LLMs struggle with complex ciphers but handle simple substitutions. \citet{lee2024decipherment} show that training language models with explicit cipher/decipherment objectives improves multilingual transfer, suggesting that cipher-awareness can be learned.

\para{In-context learning of bijective mappings.}
Most directly relevant to our work, \citet{fang2025iclciphers} introduce a framework using bijective substitution ciphers to disentangle task learning from task retrieval in in-context learning (ICL). They show that bijective ciphers consistently outperform non-bijective baselines across models and tasks (\eg +4.8\% on SST-2 for Llama-3.1-8B), with the gap widening as more demonstrations are provided. Crucially, their logit lens analysis shows that deeper layers progressively assign higher probability to the substituted token over the original, providing evidence that the residual stream can encode cipher mappings---but only when demonstrations are available.

Unlike \citet{fang2025iclciphers}, who study the \emph{with-demonstrations} setting and focus on behavioral metrics, we study the \emph{zero-shot} setting and analyze the full geometry of the residual stream. We show that without demonstrations, the model fails to recover, and that apparent representational similarity at later layers is an architectural artifact rather than genuine cipher decoding.
